% =============================================================================
%  CHAPTER 1 -- WHAT STATISTICAL ARBITRAGE ACTUALLY IS
% =============================================================================
\section{What statistical arbitrage actually is}
\label{ch:whatissaa}

\subsection{The idea, with no mathematics at all}
\label{sec:saa:idea}

Take two companies that do nearly the same thing. In this project the cleanest example
is Adobe (\code{ADBE}) and its application-software peers Salesforce (\code{CRM}),
Autodesk (\code{ADSK}), and Intuit (\code{INTU}). All four sell software to businesses,
all four are in the same GICS sector, all four are held by the same style of investor,
and all four react to the same news: interest rates, enterprise IT budgets, cloud
spending, the general appetite for growth stocks. Over years, their prices move
together --- not identically, but together.

Now suppose that on a particular day Adobe falls 3\% while its peers are flat. Nothing
Adobe-specific happened: no earnings miss, no downgrade, no lawsuit. The likely
explanation is a temporary, mechanical pressure --- a large fund redeeming and selling
its most liquid name, an index rebalance, a block trade printed at an unfortunate
price. If that is the explanation, the price gap is \emph{temporary}: it should close.

The trade writes itself. Sell Adobe short, buy the peers, and wait for the gap to
close. If Adobe recovers, the short loses but the peers (which you own) probably also
rise, and the \emph{difference} still narrows. If the whole sector falls, the short
makes money and the peers lose --- again the difference is what matters. The position
is close to indifferent to the direction of the market: it bets only on the
\emph{relationship} reverting.

That indifference is the whole point. A bet on ``the market goes up'' is a bet you
cannot win reliably; a bet on ``Adobe and Autodesk return to their normal relative
pricing'' is a much narrower, much more testable claim. Statistical arbitrage is the
industrialised version of this: find relationships that are statistically stable,
measure how far apart the names currently are, trade the gap when it is unusually
large, and hold a diversified book of such trades so that no single relationship has to
be right.

\begin{intuition}
Statistical arbitrage is not ``arbitrage'' in the textbook sense. Textbook arbitrage is
riskless: buy for \$10 here, sell for \$10.05 there, simultaneously, guaranteed profit.
Statistical arbitrage is a \emph{statistical} bet: the gap has closed often enough in
the past that we expect it to close again, on average, more often than not. It can and
does lose money --- sometimes for months --- when a relationship permanently breaks
(Adobe stops being like Autodesk because one of them changes business model). The word
``arbitrage'' survives only because the positions are hedged, not because they are
riskless.
\end{intuition}

\subsection{From pairs to baskets: why more than two legs}
\label{sec:saa:baskets}

The two-name version (``pairs trading'') is the oldest and simplest form
\cite{gatev2006pairs}. It has a well-known weakness: with only one hedge leg, all the
idiosyncratic risk of that one hedge name sits in your position. If Autodesk has an
earnings surprise, your ``market-neutral'' trade takes a direct hit that has nothing to
do with the mispricing you were betting on.

A \emph{basket} hedge fixes this by spreading the hedge over several related names. Now
the question becomes quantitative: how much of each peer should I hold against one unit
of the target? That is a regression question, and regression is the mathematical core
of this entire project:
\begin{itemize}
  \item \textbf{Which names} belong in the basket? (selection: lasso, elastic net,
    principal components --- Chapter~\ref{ch:regularization})
  \item \textbf{How much} of each? (estimation: OLS, ridge --- Chapter~\ref{ch:ols})
  \item \textbf{Do the amounts change over time?} (adaptive estimation: rolling windows,
    Kalman filter --- Chapter~\ref{ch:kalman})
  \item \textbf{Is the leftover gap stationary}, i.e.\ does it actually revert, or does
    it drift apart forever? (unit-root and cointegration tests ---
    Chapter~\ref{ch:timeseries})
  \item \textbf{How fast} does it revert, and is that fast enough to beat trading costs?
    (Ornstein--Uhlenbeck half-life --- Chapter~\ref{ch:ou})
\end{itemize}
Those five questions are the skeleton of the whole report, and they are asked in that
order in Part~II.

The systematic, factor-based version of the same idea --- regress each name on a set of
risk factors (sector ETFs, PCA factors) and trade the residual --- is
\cite{avellaneda2010statistical}, the central reference for basket construction and
residual dynamics. Our implementation is closer to the ``economic basket'' variant: the
basket is chosen for economic coherence (same sub-industry) and the weights are
estimated by regression, rather than being read off a global factor model.

\subsection{Market neutrality, precisely}
\label{sec:saa:neutrality}

``Market neutral'' is a marketing phrase unless it is defined. There are three distinct
levels, and this project enforces the first two and measures the third.

\paragraph{(1) Dollar neutrality.} The position's long notional equals its short
notional. With a target $T$ and hedge weights $\beta_1,\dots,\beta_n$ on basket names,
we short \$1 of the target and buy $\beta_i$ dollars of name $i$, then rescale the
whole book by the \emph{grossing factor}
\begin{equation}
  g \;=\; 1 + \sum_{i=1}^{n}\lvert \beta_i\rvert ,
  \qquad
  \text{positions} = \frac{1}{g}\Bigl(-1,\ \beta_1,\dots,\beta_n\Bigr),
  \label{eq:grossing}
\end{equation}
so that gross notional is exactly \$1 and net dollar exposure is
$\tfrac{1}{g}\bigl(-1+\sum_i\beta_i\bigr)$, which is zero whenever $\sum_i\beta_i=1$.
Dollar neutrality kills the first-order effect of ``everything goes up''.

\paragraph{(2) Beta neutrality.} Dollar neutrality is not enough: if the target is
twice as volatile as the basket, a dollar-neutral book still has net market exposure.
Beta neutrality requires the position's market beta to vanish,
\begin{equation}
  \beta_{\text{book}} \;=\; \sum_{i} w_i\,\beta_{i,\text{mkt}} \;=\; 0 .
  \label{eq:betaneutral}
\end{equation}
Regressing the target on economically related names does most of this work implicitly:
the fitted values $\hat y_t = \hat\alpha_t + x_t^{\top}\hat\beta_t$ absorb the common
(sector/market) component, so the residual is by construction orthogonal to the
regressors --- \emph{in sample}. Out of sample, and with time-varying betas, residual
orthogonality degrades; that is exactly what the diagnostics of
Chapter~\ref{ch:timeseries} and the portfolio risk measurements of
Phase~\ref{ch:phase7} quantify.

\paragraph{(3) Factor neutrality.} Real funds neutralize a whole risk model (market,
size, value, momentum, sector, volatility). This project does not: it has one basket
per strategy and reports the aggregate book's realized exposure. Stated as a non-goal
in Chapter~\ref{ch:nongoals}.

\begin{remark}[Honesty]
A book that is dollar-neutral and roughly beta-neutral is \emph{not} riskless. It still
carries: (i) \emph{relative-value risk} --- the spread can widen before it narrows, and
can widen permanently; (ii) \emph{liquidity risk} --- the short leg must be borrowable,
and borrow costs rise exactly when the trade is crowded; (iii) \emph{tail risk} --- in
the March 2020 dislocation, correlations went to one and relative-value books lost
money on both legs' financing. The empirical chapters report drawdowns and tail
measures, not just averages, for this reason.
\end{remark}

\subsection{Anatomy of one trade, in formulas}
\label{sec:saa:anatomy}

This subsection is a preview; every object here is derived properly in Part~I and
implemented in Part~II. The pipeline has five stages, and the C++ engine
(Chapter~\ref{ch:phase3}) has one module per stage.

\paragraph{Stage 1 --- Estimate the relationship.} At bar $t$, using only bars up to
$t-1$, fit the target's log price on the basket's log prices over a trailing window of
$W$ bars:
\begin{equation}
  y_{s} \;=\; \alpha_t + \sum_{i=1}^{n}\beta_{i,t}\,x_{i,s} + \eps_{s},
  \qquad s = t-W,\dots,t-1 .
  \label{eq:saa:regression}
\end{equation}
The subscript on $\alpha_t,\beta_{i,t}$ reminds us that these are re-estimated as the
window rolls; that time-variation is a feature, not a nuisance
(Chapter~\ref{ch:kalman}).

\paragraph{Stage 2 --- Form the spread.} Apply the just-estimated hedge to bar $t$:
\begin{equation}
  S_t \;=\; y_t - \hat\alpha_t - \sum_{i=1}^{n}\hat\beta_{i,t}\,x_{i,t}.
  \label{eq:saa:spread}
\end{equation}
$S_t$ is the mispricing candidate: how far the target is from where the basket says it
should be.

\paragraph{Stage 3 --- Standardize.} Raw spread units are meaningless across baskets and
across volatility regimes, so convert to a $z$-score using the window's own residual
statistics:
\begin{equation}
  z_t \;=\; \frac{S_t - \hat\mu_{S,t}}{\hat\sigma_{S,t}},
  \qquad
  \hat\mu_{S,t}=\frac{1}{W}\sum_{s=t-W}^{t-1}S_s,\quad
  \hat\sigma^2_{S,t}=\frac{1}{W-1}\sum_{s=t-W}^{t-1}(S_s-\hat\mu_{S,t})^2 .
  \label{eq:saa:z}
\end{equation}
Now ``$z_t=+2$'' means ``two standard deviations rich relative to its own recent
history'', comparable across names and years.

\paragraph{Stage 4 --- Decide.} A threshold policy with an entry band and an exit band:
\begin{equation}
  \text{position}_{t+1} =
  \begin{cases}
    -1 & \text{if } z_t \ge z_{\text{entry}} \quad(\text{target rich: short target, long basket})\\
    +1 & \text{if } z_t \le -z_{\text{entry}} \quad(\text{target cheap: long target, short basket})\\
    \ \ 0 & \text{if } \lvert z_t\rvert \le z_{\text{exit}} \quad(\text{converged: flat})\\
    \text{unchanged} & \text{otherwise (hold inside the bands)}.
  \end{cases}
  \label{eq:saa:policy}
\end{equation}
The hysteresis (two bands instead of one) is what keeps turnover --- and therefore cost
--- down: the book does not flip on every wiggle through zero. In this project the
baseline is $z_{\text{entry}}=2.0$, $z_{\text{exit}}=0.5$, and
Chapter~\ref{ch:phase7} maps the whole $(W, z_{\text{entry}})$ surface.

\paragraph{Stage 5 --- Execute and pay.} The decision made with data through close($t$)
is executed on bar $t+1$ --- never on bar $t$. Costs are charged on traded notional:
proportional cost $c$ (baseline 10\,bps one-way, swept 0--50\,bps) plus a square-root
market-impact term (Chapter~\ref{ch:backtestmath}). Gross and net P\&L are always both
reported.

\begin{intuition}
Read the five stages as a funnel. Stage 1--2 decide \emph{what} the signal is; stage 3
puts it in comparable units; stage 4 decides \emph{when} to act; stage 5 decides what it
\emph{costs} to act. Almost every way this kind of strategy lies to you is a bug in one
of these five stages: a non-causal stage 1 (lookahead), a stage 3 that uses full-sample
$\hat\sigma_S$ (subtle lookahead), a stage 4 tuned on the test period (data snooping),
or a stage 5 with no costs (the classic fantasy backtest). Each stage therefore gets
its own unit test in this repository.
\end{intuition}

\subsection{Where could the edge come from?}
\label{sec:saa:economics}

A profitable systematic strategy needs a reason to exist --- either someone is being
compensated for bearing a risk you are happy to bear, or someone else is making a
mistake you are equipped to exploit, or a friction prevents others from competing the
profit away. For basket stat-arb the candidate explanations are:

\begin{enumerate}
  \item \textbf{Temporary liquidity / demand pressure.} Large uninformed order flow
    (redemptions, index events, block trades) moves a single name relative to its
    economic peers. Risk-arbitrageurs with balance-sheet capacity earn a premium for
    absorbing it. This is a \emph{risk premium}: the compensation is for holding an
    position that can get worse before it gets better, and for the small probability
    that the relationship never reverts.
  \item \textbf{Limited attention / slow information diffusion.} News about a sector
    propagates across names with a lag; the laggards are predictable from the leaders.
    This is a behavioural explanation and it decays as more capital watches for it.
  \item \textbf{Frictions.} Short-sale costs, borrow constraints, and margin
    requirements keep mispricings alive that would otherwise be trivially arbitraged.
    The same frictions apply to you: your net return is the gross mispricing minus the
    cost of the borrow and the cost of the round trip.
  \item \textbf{It does not exist (any more).} The honest default. Pairs trading
    profitability has been documented to decay \cite{dofaff2010does}, and the most
    crowded relative-value signals are the first to be competed away.
\end{enumerate}

This project cannot distinguish between explanations 1--3, and says so. What it
\emph{can} do is measure the size of the effect after costs, and ask whether that size
is distinguishable from zero once the number of things searched for is taken into
account. That is Chapters~\ref{ch:backtestmath}--\ref{ch:multipletesting}, and the
answer here is ``not distinguishable'' --- which is a legitimate, publishable, and
useful result.

\subsection{The four ways this kind of project lies to you}
\label{sec:saa:lies}

Everything in the design of this repository is a response to one of these four.

\begin{enumerate}[label=\textbf{Lie \arabic*.}]
  \item \textbf{Lookahead bias.} Using information at time $t$ that was not available
    until after $t$. The subtle forms are the killers: standardizing by the
    full-sample $\sigma$, choosing the basket after seeing which spread reverted best,
    filling missing prices with the next available value. \emph{Defence here:} the
    estimation window ends at $t-1$, execution is at $t+1$, one C++ unit test
    manufactures a single large mispricing and asserts the position is never active on
    the signal day, and the Kalman test asserts $\beta_{t\mid t-1}$ depends only on bars
    $<t$.
  \item \textbf{Survivorship bias.} Using a universe of names that all survived to the
    end of the sample. Companies that were acquired, delisted, or went bankrupt
    disappear from most free data sources, and their disappearance flatters every
    strategy. \emph{Defence here:} none available in-sandbox --- the panel is
    survivorship-biased and Chapter~\ref{ch:phase1} quantifies and discloses it rather
    than pretending. This is a real limitation and it biases results \emph{upward}.
  \item \textbf{Data snooping / multiple testing.} Try 44 configurations, report the
    best, call it a discovery. With 44 independent null strategies, the chance that at
    least one shows $p<0.05$ is $1-0.95^{44}\approx 89\%$. \emph{Defence here:} a full
    hypothesis-count audit, Bonferroni, White's Reality Check, Hansen's SPA, and
    pairwise Diebold--Mariano tests (Chapter~\ref{ch:multipletesting},
    Phase~\ref{ch:phase10}).
  \item \textbf{Ignoring costs.} A daily-frequency strategy that turns over 133--362
    round trips over sixteen years pays, at 10\,bps one-way, a drag that in this
    project is \emph{larger than the gross edge}. \emph{Defence here:} every result is
    reported gross and net, with a 0--50\,bps sensitivity sweep.
\end{enumerate}

\subsection{The base model, stated once and for all}
\label{sec:basemodel}

For a target asset $Y_t$ and a basket of $n$ candidate assets $X_{1,t},\dots,X_{n,t}$,
\begin{equation}
  Y_t \;=\; \alpha + \sum_{i=1}^{n}\beta_i\,X_{i,t} + \eps_t ,
  \label{eq:basemodel}
\end{equation}
where $\eps_t$ is the tradeable signal. Which $Y$ and $X$ to use (levels or logs or
returns), which $X$'s to include, how to estimate $\beta$, how often to re-estimate it,
and how to map $\eps_t$ into positions and pay for the privilege --- that is the entire
substance of the following forty-odd pages.

\subsection{What this project contributes}
\label{sec:saa:contribution}

Not a new alpha. Instead: a complete, reproducible, adversarially self-critiqued
implementation of the basket stat-arb idea, in which
\begin{itemize}
  \item the statistical layer is a tested library (\code{scripts/statkit/}, 35 unit
    tests on synthetic data where the truth is known, so the tests can actually fail);
  \item the simulation layer is a dependency-free C++ engine with per-stage latency
    histograms, allocation counting, and no-lookahead assertions (\code{src/},
    \code{test/}, \code{bench/});
  \item the two layers are cross-checked against each other (brute-force vs.\
    incremental regression to $10^{-8}$; Python vs.\ C++ walk-forward hedges);
  \item every headline number is accompanied by the correction that could destroy it;
  \item the report is \emph{built} from \code{results/} by
    \code{report/build_report.py}, so the numbers cannot drift from the code.
\end{itemize}

The value of a negative result delivered this way is that it is \emph{falsifiable and
re-runnable}: you can change the universe, the cost, or the penalty and see exactly
which conclusions move.
